% !TeX program = lualatex
% Agregátor reálných otázek SZZ pro okruh I1. Jednotlivé otázky jsou v souborech q_*.tex.
% Sestavení: ./build.sh 02_spolecna_informatika/I1_automaty_jazyky/priklady   (nebo cd .../priklady && latexmk)
\documentclass[11pt,a4paper]{article}
\usepackage{szzpriklady}
% Build běží z adresáře priklady/ (CWD), obrázky jsou v src/fig/.
\graphicspath{{src/fig/}{fig/}}

\title{Příklady otázek -- Automaty a jazyky\\[0.3em]\large Reálné otázky ze SZZ 2023--2026 (I1)}
\author{Jakub Klesa}
\date{Pokryté předměty: NTIN071 (Automaty a gramatiky)}

\begin{document}
\maketitle
\tableofcontents
\bigskip

\begin{poznamka}
Tento dokument obsahuje reálné otázky z minulých termínů SZZ (2023--2026) zařazené do
okruhu I1. U otázek s obrázkem je grafika oříznutá z originálního PDF (viz odkaz na
zdroj). Text, číselná data a tabulky jsou přepsané věrně.
Zadání i oficiální náčrty řešení pocházejí ze sad zveřejněných na webu MFF UK:\par
{\raggedright\url{https://www.mff.cuni.cz/cs/studenti/bakalarske-studium/statni-zaverecne-zkousky/bakalarske-statni-zkousky-studijniho-programu-informatika}\par}
Náčrty označené jako doplněné jsou vlastní, oficiální sada je k~dané otázce neobsahuje.
\end{poznamka}
\bigskip

% === Otázky (chronologicky). Doplň \input řádky a soubory q_*.tex do src/. ===
% Příklad: \input{q_2024-leto-28_bayesovske-site}
% ZACATEK-OTAZEK
\input{src/q_2023-leto-02_cfg-na-zasobnikovy-automat}
\input{src/q_2023-podzim-02_zasobnikovy-automat}
\input{src/q_2024-jaro-01_reprezentace-bezkontextoveho-jazyka}
\input{src/q_2024-podzim-01_doplnek-regularniho-vyrazu}
\input{src/q_2025-jaro-01_prunik-bezkontextoveho-regularniho}
\input{src/q_2025-leto-01_deterministicke-konecne-automaty}
\input{src/q_2026-jaro-01_porovnani-modularnich-hashu}
% KONEC-OTAZEK

\end{document}
